\documentclass[12pt, a4paper]{article}

% ── Packages ──────────────────────────────────────────────────────────────────
\usepackage[T1]{fontenc}
\usepackage[utf8]{inputenc}
\usepackage{lmodern}
\usepackage[margin=2.5cm]{geometry}
\usepackage{amsmath, amssymb, amsthm}
\usepackage{mathtools}
\usepackage{bm}
\usepackage{bbm}
\usepackage{microtype}
\usepackage{graphicx}
\usepackage{booktabs}
\usepackage{array}
\usepackage{multirow}
\usepackage{xcolor}
\usepackage{hyperref}
\usepackage{cleveref}
\usepackage{enumitem}
\usepackage{fancyhdr}
\usepackage{titlesec}
\usepackage{abstract}
\usepackage{setspace}
\usepackage{natbib}
\usepackage{listings}
\usepackage[ruled, vlined, linesnumbered]{algorithm2e}
\usepackage{tcolorbox}
\tcbuselibrary{skins, breakable}

% ── Hyperref setup ────────────────────────────────────────────────────────────
\hypersetup{
  colorlinks  = true,
  linkcolor   = blue!70!black,
  citecolor   = blue!70!black,
  urlcolor    = blue!70!black,
  pdftitle    = {Detecting Structural Breaks via the Adaptive Fused Lasso},
  pdfauthor   = {Yousef Kaddoura},
}

% ── Page style ────────────────────────────────────────────────────────────────
\pagestyle{fancy}
\fancyhf{}
\rhead{\small Kaddoura --- Fused Lasso Break Detection}
\lhead{\small Technical Documentation v1.0}
\cfoot{\thepage}
\renewcommand{\headrulewidth}{0.4pt}

% ── Title formatting ──────────────────────────────────────────────────────────
\titleformat{\section}{\large\bfseries}{\thesection.}{0.75em}{}[\titlerule]
\titleformat{\subsection}{\normalsize\bfseries}{\thesubsection.}{0.5em}{}
\titleformat{\subsubsection}{\normalsize\itshape}{\thesubsubsection.}{0.5em}{}

% ── Theorem environments ──────────────────────────────────────────────────────
\theoremstyle{plain}
\newtheorem{assumption}{Assumption}[section]
\newtheorem{theorem}{Theorem}[section]
\newtheorem{proposition}{Proposition}[section]
\newtheorem{lemma}{Lemma}[section]
\newtheorem{corollary}{Corollary}[section]

\theoremstyle{definition}
\newtheorem{definition}{Definition}[section]
\newtheorem{example}{Example}[section]

\theoremstyle{remark}
\newtheorem{remark}{Remark}[section]

% ── Custom notation ───────────────────────────────────────────────────────────
\newcommand{\R}{\mathbb{R}}
\newcommand{\E}{\mathbb{E}}
\newcommand{\bB}{\mathcal{B}}
\newcommand{\N}{\mathcal{N}}
\newcommand{\bbeta}{\bm{\beta}}
\newcommand{\blambda}{\bm{\lambda}}
\newcommand{\bF}{\mathbf{F}}
\newcommand{\bX}{\tilde{\mathbf{X}}}
\newcommand{\by}{\tilde{\mathbf{y}}}
\newcommand{\bu}{\tilde{\mathbf{u}}}
\newcommand{\norm}[1]{\left\| #1 \right\|}
\newcommand{\abs}[1]{\left| #1 \right|}
\newcommand{\floor}[1]{\left\lfloor #1 \right\rfloor}
\newcommand{\plim}{\operatorname{plim}}

% ── Code listing style ────────────────────────────────────────────────────────
\definecolor{codebackground}{rgb}{0.97,0.97,0.97}
\definecolor{codekeyword}{rgb}{0.0, 0.3, 0.7}
\definecolor{codestring}{rgb}{0.5, 0.0, 0.0}
\definecolor{codecomment}{rgb}{0.3, 0.5, 0.3}
\lstdefinestyle{pythonstyle}{
  backgroundcolor  = \color{codebackground},
  basicstyle       = \ttfamily\footnotesize,
  keywordstyle     = \color{codekeyword}\bfseries,
  stringstyle      = \color{codestring},
  commentstyle     = \color{codecomment}\itshape,
  numberstyle      = \tiny\color{gray},
  numbers          = left,
  numbersep        = 6pt,
  showstringspaces = false,
  breaklines       = true,
  frame            = single,
  framesep         = 3pt,
  rulecolor        = \color{gray!40},
  language         = Python,
  tabsize          = 4,
  captionpos       = b,
}
\lstset{style=pythonstyle}

% ── Info box ──────────────────────────────────────────────────────────────────
\newtcolorbox{infobox}[1][]{
  enhanced, breakable,
  colback  = blue!5!white,
  colframe = blue!60!black,
  fonttitle= \bfseries,
  title    = {#1},
  left=5pt, right=5pt, top=4pt, bottom=4pt,
}

% ─────────────────────────────────────────────────────────────────────────────
\begin{document}
% ─────────────────────────────────────────────────────────────────────────────

% ── Title page ───────────────────────────────────────────────────────────────
\begin{titlepage}
\centering
\vspace*{1.5cm}

{\LARGE\bfseries Detecting Structural Breaks via the\\[0.4em]
Adaptive Fused Lasso\par}

\vspace{0.5cm}
{\large Technical Documentation and Replication Guide\par}

\vspace{1.5cm}

\rule{0.6\textwidth}{0.8pt}

\vspace{1cm}

{\large \textbf{Yousef Kaddoura}\par}
\vspace{0.3cm}
{\normalsize
Department of Economics\\
\vspace{0.15cm}
\href{mailto:yousef.kaddoura@example.com}{\texttt{yousef.kaddoura@example.com}}
\par}

\vspace{1cm}

\rule{0.6\textwidth}{0.8pt}

\vspace{1cm}

{\normalsize
Based on:\\[0.3em]
\textit{Estimation of panel data models with random interactive effects\\
and multiple structural breaks when $T$ is fixed}\\[0.3em]
Kaddoura, Y.\ and Westerlund, J.\ (2023)\\
\textbf{Journal of Business \& Economic Statistics}, 41, pp.\ 778--790
\par}

\vspace{2cm}

{\normalsize Repository version: \texttt{v1.0.0} \hfill \today\par}

\end{titlepage}

% ── Table of contents ─────────────────────────────────────────────────────────
\tableofcontents
\newpage

% ─────────────────────────────────────────────────────────────────────────────
\begin{onehalfspacing}
% ─────────────────────────────────────────────────────────────────────────────

% ═════════════════════════════════════════════════════════════════════════════
\section{Introduction}
% ═════════════════════════════════════════════════════════════════════════════

A central concern in empirical macroeconomics and finance is whether the
relationship between variables is stable over time.
When this relationship undergoes an abrupt shift --- a \emph{structural break}
--- standard estimators that assume parameter constancy yield inconsistent
estimates and invalid inference.

The contribution of \citet{kaddoura2023estimation} lies in developing a
unified, computationally tractable estimator that simultaneously:
\begin{enumerate}[noitemsep]
  \item eliminates latent interactive fixed effects through cross-sectional
        demeaning;
  \item detects an unknown number of structural breaks without prior
        specification;
  \item dates those breaks consistently, even when the time dimension $T$
        is small and fixed.
\end{enumerate}
The key ingredient is the \textbf{adaptive fused lasso}, which penalises
consecutive differences in the coefficient vector rather than the coefficients
themselves, encouraging piecewise-constant solutions.

This document provides a self-contained technical exposition of the model,
estimation procedure, and Monte Carlo design, alongside a guide to the
companion Python codebase.

% ═════════════════════════════════════════════════════════════════════════════
\section{Model and Notation}
\label{sec:model}
% ═════════════════════════════════════════════════════════════════════════════

\subsection{Panel Data Model}
\label{sub:panel}

Consider the panel regression
\begin{equation}
  \label{eq:model}
  y_{it} = x_{it}'\,\beta_t + u_{it},
  \qquad i = 1,\ldots,N, \quad t = 1,\ldots,T,
\end{equation}
where:
\begin{itemize}[noitemsep]
  \item $y_{it} \in \R$ is the dependent variable for unit $i$ at time $t$;
  \item $x_{it} \in \R^p$ is a vector of regressors;
  \item $\bbeta_t \in \R^p$ is the slope coefficient vector, allowed to vary
        across $t$;
  \item $u_{it}$ is the composite error.
\end{itemize}

\subsection{Interactive Fixed Effects}
\label{sub:ife}

The error is decomposed as
\begin{equation}
  \label{eq:error}
  u_{it} = \lambda_i' F_t + \varepsilon_{it},
\end{equation}
where $F_t \in \R^r$ are unobserved common factors, $\lambda_i \in \R^r$ are
unit-specific factor loadings, and $\varepsilon_{it}$ is an idiosyncratic
shock.  The term $\lambda_i' F_t$ represents \emph{interactive fixed effects},
generalising the additive two-way fixed effect model.

\begin{remark}
  Standard within-group transformations (fixed effect demeaning over $i$ or
  $t$) do not eliminate interactive fixed effects because the cross product
  $\lambda_i' F_t$ cannot be separated by such linear projections.
\end{remark}

\subsection{Cross-Sectional Demeaning}
\label{sub:demean}

Following the literature on factor-augmented regressions, we eliminate the
common component by cross-sectionally demeaning each variable.  Define
\begin{equation}
  \label{eq:demean}
  \tilde{z}_{it} \coloneqq z_{it} - \bar{z}_t, \qquad
  \bar{z}_t \coloneqq \frac{1}{N}\sum_{i=1}^N z_{it},
\end{equation}
for $z \in \{y, x_1, \ldots, x_p\}$.  Under standard assumptions on the
factor loadings (see \cref{sec:assumptions}), the demeaned model satisfies
\begin{equation}
  \label{eq:demeaned}
  \tilde{y}_{it} = \tilde{x}_{it}'\,\bbeta_t + \tilde{u}_{it},
\end{equation}
where $\tilde{u}_{it}$ has mean zero across $i$ and the common factor
component is asymptotically negligible as $N \to \infty$.

\subsection{Structural Breaks}
\label{sub:breaks}

The coefficient vector $\bbeta_t$ follows a piecewise-constant path with
$m$ unknown breakpoints $t_1^* < t_2^* < \cdots < t_m^*$:
\begin{equation}
  \label{eq:breaks}
  \bbeta_t =
  \begin{cases}
    \bbeta^{(0)} & 1 \le t \le t_1^*, \\
    \bbeta^{(1)} & t_1^*+1 \le t \le t_2^*, \\
    \quad \vdots \\
    \bbeta^{(m)} & t_m^*+1 \le t \le T.
  \end{cases}
\end{equation}
The number of breaks $m$ and their locations $\{t_j^*\}$ are treated as
\emph{unknown} and estimated jointly with the coefficients.

% ═════════════════════════════════════════════════════════════════════════════
\section{Estimation Procedure}
\label{sec:estimation}
% ═════════════════════════════════════════════════════════════════════════════

\subsection{First Stage: Time-Varying OLS}
\label{sub:ols}

As a first step, compute the unrestricted period-by-period OLS estimator:
\begin{equation}
  \label{eq:ols}
  \hat{\bbeta}_t^{\text{OLS}} =
    \underset{\bbeta}{\arg\min}\;
    \frac{1}{N} \norm{\tilde{y}_t - \tilde{X}_t \bbeta}^2,
    \qquad t = 1,\ldots,T,
\end{equation}
where $\tilde{y}_t = (\tilde{y}_{1t},\ldots,\tilde{y}_{Nt})' \in \R^N$ and
$\tilde{X}_t \in \R^{N \times p}$.  This is equivalent to the analytical
solution $\hat{\bbeta}_t^{\text{OLS}} = (\tilde{X}_t'\tilde{X}_t)^{-1}\tilde{X}_t'\tilde{y}_t$
whenever $\tilde{X}_t'\tilde{X}_t$ is invertible.

The OLS estimates do not impose any smoothness or piecewise-constancy, so
they serve exclusively to construct the adaptive weights in the second stage.

\subsection{Adaptive Weights}
\label{sub:weights}

For each $t \ge 2$, define the adaptive weight
\begin{equation}
  \label{eq:weights}
  w_t = \norm{\hat{\bbeta}_t^{\text{OLS}} - \hat{\bbeta}_{t-1}^{\text{OLS}}}_2^{-2}.
\end{equation}
The weight $w_t$ is \emph{small} (penalty is light) when the OLS path already
shows a large jump at $t$, allowing the fused lasso to place a break there
easily. Conversely, $w_t$ is \emph{large} (penalty is heavy) when the OLS
path is flat at $t$, discouraging spurious breaks.

\begin{remark}[Numerical stability]
  When $\norm{\hat{\bbeta}_t^{\text{OLS}} - \hat{\bbeta}_{t-1}^{\text{OLS}}}_2 < \delta$
  for some small $\delta > 0$, the weight is capped at $\delta^{-2}$ to avoid
  division by zero. In the code, $\delta = 10^{-12}$.
\end{remark}

\subsection{Second Stage: Adaptive Fused Lasso}
\label{sub:fgls}

The fused lasso estimator solves the penalised least-squares problem:
\begin{equation}
  \label{eq:fused_lasso}
  \hat{\bB}_T(\lambda) =
    \underset{\bB_T}{\arg\min} \;
    \underbrace{\frac{1}{N}\sum_{i=1}^N \sum_{t=1}^T
      \bigl(\tilde{y}_{it} - \tilde{x}_{it}'\,\bbeta_t\bigr)^2}_{\text{fit}}
    + \;
    \underbrace{\lambda \sum_{t=2}^T w_t
      \norm{\bbeta_t - \bbeta_{t-1}}_2}_{\text{fused penalty}},
\end{equation}
where $\bB_T = (\bbeta_1, \ldots, \bbeta_T) \in \R^{p \times T}$ collects all
coefficient vectors and $\lambda > 0$ is the regularisation parameter.

\begin{infobox}[Interpretation of the fused penalty]
  The penalty $\lambda w_t \norm{\bbeta_t - \bbeta_{t-1}}_2$ shrinks
  \emph{differences} between consecutive coefficient vectors toward zero.
  When the penalty is large enough, two adjacent vectors become identical ---
  the lasso sets $\hat{\bbeta}_t = \hat{\bbeta}_{t-1}$ exactly --- which
  corresponds to \emph{no break} between $t-1$ and $t$. The locations where
  the solution is not exactly zero identify the estimated break dates.
\end{infobox}

Problem \eqref{eq:fused_lasso} is a \emph{convex} second-order cone program
and is solved globally by the CVXPY/CLARABEL solver.

\subsection{Break Detection}
\label{sub:detection}

Given the solution $\hat{\bB}_T(\lambda)$, the estimated number of breaks is
\begin{equation}
  \label{eq:mhat}
  \hat{m}(\lambda) = \sum_{t=2}^T
    \mathbf{1}\!\left\{
      \norm{\hat{\bbeta}_t(\lambda) - \hat{\bbeta}_{t-1}(\lambda)}_2
      > \tau_{\text{tol}}
    \right\},
\end{equation}
with the tolerance $\tau_{\text{tol}} = 10^{-3}$ by default.  The estimated
break dates are the values of $t$ for which the indicator equals one.

% ═════════════════════════════════════════════════════════════════════════════
\section{Tuning Parameter Selection}
\label{sec:tuning}
% ═════════════════════════════════════════════════════════════════════════════

The regularisation parameter $\lambda$ controls the trade-off between model
fit and model complexity (number of breaks).  We select $\lambda$ by
minimising an information criterion:

\begin{equation}
  \label{eq:ic}
  \text{IC}(\lambda) =
    \underbrace{
      \frac{1}{NT}\sum_{t=1}^T
      \norm{\tilde{y}_t - \tilde{X}_t\,\hat{\bbeta}_t(\lambda)}_2^2
    }_{\text{in-sample fit}}
    + \;
    \underbrace{
      \frac{\log N}{\sqrt{N}} \cdot p \cdot \bigl(\hat{m}(\lambda)+1\bigr)
    }_{\text{complexity penalty}}.
\end{equation}

The factor $\log(N)/\sqrt{N}$ ensures that the penalty grows with the
cross-sectional dimension $N$ but at a rate slower than $1/\sqrt{N}$, which
is the standard rate of convergence of the OLS estimator.  The optimal
parameter is
\begin{equation}
  \label{eq:lamstar}
  \lambda^* = \underset{\lambda \in \Lambda}{\arg\min}\; \text{IC}(\lambda),
\end{equation}
where $\Lambda = \{10^{-3}, \ldots, 10^{3}\}$ is a log-spaced grid of 50
candidate values.

\begin{algorithm}[H]
  \caption{Adaptive Fused Lasso for Break Detection}
  \KwIn{Demeaned data $(\tilde{Y}, \tilde{X})$, grid $\Lambda$,
        tolerance $\tau_{\text{tol}}$}
  \KwOut{$\hat{\bB}_T$, $\hat{m}$, $\lambda^*$}
  \BlankLine
  \tcp{Stage 1: OLS initialisation}
  Compute $\hat{\bB}^{\text{OLS}}$ by solving \eqref{eq:ols} for each $t$\;
  \BlankLine
  \tcp{Stage 2: IC minimisation over $\Lambda$}
  \ForEach{$\lambda \in \Lambda$}{
    Compute adaptive weights $w_t$ via \eqref{eq:weights}\;
    Solve fused lasso \eqref{eq:fused_lasso} to obtain $\hat{\bB}_T(\lambda)$\;
    Count breaks $\hat{m}(\lambda)$ via \eqref{eq:mhat}\;
    Evaluate $\text{IC}(\lambda)$ via \eqref{eq:ic}\;
  }
  Set $\lambda^* \leftarrow \arg\min_{\lambda \in \Lambda}\,\text{IC}(\lambda)$\;
  \BlankLine
  \tcp{Stage 3: Final estimation}
  Solve \eqref{eq:fused_lasso} at $\lambda^*$ to obtain
  $\hat{\bB}_T \coloneqq \hat{\bB}_T(\lambda^*)$\;
  Set $\hat{m} \leftarrow \hat{m}(\lambda^*)$\;
  \Return $\hat{\bB}_T,\; \hat{m},\; \lambda^*$\;
\end{algorithm}

% ═════════════════════════════════════════════════════════════════════════════
\section{Assumptions}
\label{sec:assumptions}
% ═════════════════════════════════════════════════════════════════════════════

We state the high-level assumptions required for consistency and break-date
localisation.  Full primitive conditions are in \citet{kaddoura2023estimation}.

\begin{assumption}[Factor Loadings]
  \label{ass:loadings}
  $\{\lambda_i\}_{i=1}^N$ are i.i.d.\ with $\E[\lambda_i] = \mu_\lambda$ and
  $\E[\norm{\lambda_i}^4] < \infty$.  Moreover,
  $N^{-1}\sum_i \lambda_i \lambda_i' \xrightarrow{p} \Sigma_\lambda$ as
  $N \to \infty$ with $\Sigma_\lambda$ positive definite.
\end{assumption}

\begin{assumption}[Errors]
  \label{ass:errors}
  $\varepsilon_{it}$ is weakly cross-sectionally and temporally dependent with
  $\E[\varepsilon_{it}] = 0$, $\E[\varepsilon_{it}^2] < \infty$, and
  $\sup_{i,t}\E[\abs{\varepsilon_{it}}^{4+\delta}] < \infty$ for some $\delta > 0$.
\end{assumption}

\begin{assumption}[Regressors]
  \label{ass:regressors}
  $x_{it}$ may be correlated with the common factors $F_t$ (through factor
  loadings $\lambda_{ik}$) but is otherwise weakly dependent.
\end{assumption}

\begin{assumption}[Break Magnitude]
  \label{ass:breaks}
  Each break satisfies $\norm{\bbeta^{(j)} - \bbeta^{(j-1)}}_2 \ge c_N > 0$
  for some sequence $c_N \to \infty$ as $N \to \infty$.
\end{assumption}

\begin{assumption}[Number of Breaks]
  \label{ass:number}
  $m$ is fixed and finite.
\end{assumption}

Under Assumptions \ref{ass:loadings}--\ref{ass:number}, \citet{kaddoura2023estimation}
establish that $\hat{m}(\lambda^*) \xrightarrow{p} m$ and that each break date
is consistently estimated.

% ═════════════════════════════════════════════════════════════════════════════
\section{Data Generating Processes}
\label{sec:dgp}
% ═════════════════════════════════════════════════════════════════════════════

The Monte Carlo study employs three DGP families of increasing complexity.

\subsection{DATA1: i.i.d.\ Errors}

\begin{align}
  y_{it} &= x_{it}'\,\bbeta_t + \varepsilon_{it},\\
  x_{it} &\sim \mathcal{N}(0, I_p), \quad \varepsilon_{it} \sim \mathcal{N}(0,1),
\end{align}
all independently across $i$ and $t$.  There is no interactive fixed effect.

\subsection{DATA2: Factor-Structured Errors}

\begin{align}
  u_{it} &= \lambda_i' F_t + \varepsilon_{it},\\
  \lambda_i &\sim \mathcal{N}(\mathbf{1}_r, I_r), \quad
  F_t \sim \mathcal{N}(0, I_r), \quad \varepsilon_{it} \sim \mathcal{N}(0,1),
\end{align}
inducing cross-sectional dependence through the factor structure, while
regressors remain i.i.d.\ normal.

\subsection{DATA3: Main DGP (AR(1) Factors + Spatial Errors)}
\label{sub:data3}

This is the richest DGP and the one used in the main simulation study:

\paragraph{Common factors.}
\begin{equation}
  F_t = (1 - \phi)\,\mathbf{1}_r + \phi\,F_{t-1} + \eta_t, \qquad
  \eta_t \sim \mathcal{N}(0, I_r),
\end{equation}
with $F_0 = \mathbf{0}$ and persistence parameter $\phi \in (0,1)$.

\paragraph{Idiosyncratic errors.}
\begin{equation}
  \varepsilon_{it} = \phi_1\,\varepsilon_{i,t-1}
    + \xi_{it}
    + \phi_1 \sum_{j \in \mathcal{N}(i)} \xi_{jt},
\end{equation}
where $\xi_{it} \sim \mathcal{N}(0, \sigma_i^2)$ with
$\sigma_i \sim \mathrm{Uniform}(0.5, 1)$, and $\mathcal{N}(i)$ denotes the
set of up to ten nearest spatial neighbours of unit $i$ (wrapping around).

\paragraph{Regressors.}
\begin{equation}
  X_{itk} = F_t'\,\lambda_{ik} + \nu_{itk},
\end{equation}
where $\lambda_{ik} \sim \mathcal{N}(2\,\mathbf{1}_r, I_r)$ and $\nu_{itk}$
follows the same spatial-AR(1) structure as $\varepsilon_{it}$, with
persistence parameter $\pi$.

\paragraph{Composite error.}
\begin{equation}
  u_{it} = F_t'\,\lambda_i^{(y)} + \varepsilon_{it}, \qquad
  \lambda_i^{(y)} \sim \mathcal{N}(2\,\mathbf{1}_r, I_r).
\end{equation}

The full parameter vector is $(\phi, \phi_1, \pi)$.  In the main simulation,
we set $\phi = 0.8$, $\phi_1 = 0.4$, $\pi = 0.4$.

\subsection{Break Scenarios}

Each DGP class implements four break scenarios, summarised in
\cref{tab:scenarios}.

\begin{table}[htb]
  \centering
  \caption{Break scenarios across DGP methods.}
  \label{tab:scenarios}
  \begin{tabular}{lll}
    \toprule
    Method & Break configuration & $\bbeta_t$ assignment \\
    \midrule
    \texttt{DGP1} & 1 break at $t^* = \floor{T/2}$
      & $\bbeta_t = j\cdot\mathbf{1}_p$ for regime $j\in\{0,1\}$ \\
    \texttt{DGP2} & 2 breaks at $\floor{T/3}$ and $\floor{2T/3}$
      & $\bbeta_t = j\cdot\mathbf{1}_p$ for regime $j\in\{0,1,2\}$ \\
    \texttt{DGPA} & Break at every $t$
      & $\bbeta_t = t\cdot\mathbf{1}_p$ \\
    \texttt{DGPO} & No breaks
      & $\bbeta_t = \mathbf{0}$ for all $t$ \\
    \bottomrule
  \end{tabular}
\end{table}

% ═════════════════════════════════════════════════════════════════════════════
\section{Monte Carlo Design}
\label{sec:mc}
% ═════════════════════════════════════════════════════════════════════════════

\subsection{Configuration: N = 25, T = 5, m = 1}
\label{sub:N25T5m1}

The reference experiment (file \texttt{simulations/N25\_T5\_m1.py}) uses the
dimensions and parameters listed in \cref{tab:config}.

\begin{table}[htb]
  \centering
  \caption{Configuration for the N=25, T=5, m=1 Monte Carlo study.}
  \label{tab:config}
  \begin{tabular}{lll}
    \toprule
    Parameter & Value & Description \\
    \midrule
    $N$              & 25      & Cross-sectional dimension \\
    $T$              & 5       & Time dimension (fixed) \\
    $p$              & 4       & Number of regressors \\
    $m$              & 1       & Number of true structural breaks \\
    $r$              & 5       & Number of latent factors \\
    $t^*$            & 2       & True break date $= \floor{5/2}$ \\
    $\phi$           & 0.8     & Factor AR(1) persistence \\
    $\phi_1$         & 0.4     & Error spatial-temporal weight \\
    $\pi$            & 0.4     & Regressor spatial-temporal weight \\
    $S$              & 1\,000  & Monte Carlo replications \\
    $|\Lambda|$      & 50      & Grid size for $\lambda$ \\
    $\Lambda$        & $[10^{-3}, 10^{3}]$ & Log-spaced grid \\
    $\tau_\text{tol}$& $10^{-2}$ & Break detection threshold \\
    \bottomrule
  \end{tabular}
\end{table}

The true coefficient matrices are:
\begin{equation}
  \bbeta_t =
  \begin{cases}
    \mathbf{0}_p & t \in \{1, 2\}, \\
    \mathbf{1}_p & t \in \{3, 4, 5\}.
  \end{cases}
\end{equation}

\subsection{Performance Metrics}

For each simulation draw $s = 1,\ldots,S$ (excluding failed replications):
\begin{itemize}[noitemsep]
  \item $\hat{m}^{(s)}$ — estimated number of breaks.
  \item $\hat{\bB}_T^{(s)}$ — estimated coefficient matrix.
\end{itemize}

We report:

\paragraph{Average estimated break count.}
\begin{equation}
  \bar{m} = \frac{1}{S - S_{\text{fail}}} \sum_{s=1}^{S} \hat{m}^{(s)}.
\end{equation}

\paragraph{Frequency of incorrect break count.}
\begin{equation}
  \text{Freq}_m = \sum_{s=1}^{S} \mathbf{1}\!\left\{\hat{m}^{(s)} \ne m\right\}.
\end{equation}

\paragraph{Frequency of incorrect break date (conditional on correct count).}
Among replications where $\hat{m}^{(s)} = m$, count those for which the
detected break does not correspond to the true date $t^*$:
\begin{equation}
  \text{Freq}_t = \sum_{s:\,\hat{m}^{(s)}=m}
    \mathbf{1}\!\left\{
      \norm{\hat{\bbeta}_{t^*}^{(s)} - \hat{\bbeta}_{t^*-1}^{(s)}}_2
      < \tau_\text{tol}
    \right\}.
\end{equation}

\paragraph{Mean Frobenius Error (MFE).}
\begin{equation}
  \label{eq:mfe}
  \text{MFE} = \frac{1}{pT}
    \norm{\bB_T^{\text{true}} - \bar{\bB}_T}_F, \qquad
  \bar{\bB}_T = \frac{1}{S - S_{\text{fail}}}
    \sum_s \hat{\bB}_T^{(s)}.
\end{equation}

% ═════════════════════════════════════════════════════════════════════════════
\section{Codebase Reference}
\label{sec:code}
% ═════════════════════════════════════════════════════════════════════════════

\subsection{Repository Structure}

\begin{verbatim}
Detecting-Breaks-Via-the-Fused-Lasso/
│
├── src/
│   ├── dgp.py           # DATA1, DATA2, DATA3
│   ├── estimator.py     # Optimize (OLS, FGLS, NBOLS)
│   ├── ic.py            # information_criterion()
│   └── utils.py         # print_mc_summary(), plot_*()
│
├── simulations/
│   └── N25_T5_m1.py     # Main replication script
│
├── tests/
│   └── test_core.py     # pytest test suite
│
├── docs/
│   └── documentation.tex  ← this file
│
├── environment.yml
└── pyproject.toml
\end{verbatim}

\subsection{Key Classes and Functions}

\subsubsection{\texttt{src.dgp} — Data Generating Processes}

\begin{center}
\begin{tabular}{lll}
  \toprule
  Class & Constructor & Methods \\
  \midrule
  \texttt{DATA1} & \texttt{(m, T, n, p)} & \texttt{DGP1, DGP2, DGPA, DGPO} \\
  \texttt{DATA2} & \texttt{(r, m, T, n, p)} & \texttt{DGP1, DGP2, DGPA, DGPO} \\
  \texttt{DATA3} & \texttt{(r, m, T, n, p, phi, phi\_1, pi)} & \texttt{DGP1, DGP2, DGPA, DGPO} \\
  \bottomrule
\end{tabular}
\end{center}

Each \texttt{DGP1} call returns a tuple
\texttt{(X, y, beta, u, eps, F, y\_tilde, u\_tilde, X\_mean, X\_tilde)}.

\subsubsection{\texttt{src.estimator.Optimize} — Estimation Engine}

\begin{lstlisting}[caption={Using the Optimize class.}]
from src.estimator import Optimize
opt = Optimize(p=4, T=5, n=25)

# Period-by-period OLS
b_ols, status, val = opt.OLS(X_tilde, y_tilde)

# Adaptive fused lasso
b_hat, m_hat, status, val = opt.FGLS(
    X_tilde, y_tilde, b_ols, lam=0.1, break_tol=1e-3
)
\end{lstlisting}

\subsubsection{\texttt{src.ic.information\_criterion} — Tuning Selection}

\begin{lstlisting}[caption={IC-based selection of the regularisation parameter.}]
import numpy as np
from src.ic import information_criterion

lam_grid = np.logspace(-3, 3, 50)
IC_vec, m_brk, IC_min, idx, lam_star, m_star = information_criterion(
    lam_grid, y_tilde, X_tilde, p=4, T=5, n=25
)
print(f"Optimal lambda: {lam_star:.4f}, estimated breaks: {m_star}")
\end{lstlisting}

\subsubsection{\texttt{src.utils} — Reporting and Visualisation}

\begin{lstlisting}[caption={Generating diagnostic plots.}]
from src.utils import plot_ic_curve, plot_beta_path

# IC curve with detected lambda
fig1 = plot_ic_curve(
    lam_grid, IC_vec, m_brk, lam_star,
    save_path="figures/ic_curve.pdf"
)

# True vs estimated coefficients
fig2 = plot_beta_path(
    beta_true, b_hat,
    save_path="figures/beta_path.pdf"
)
\end{lstlisting}

% ═════════════════════════════════════════════════════════════════════════════
\section{Installation and Reproducibility}
\label{sec:install}
% ═════════════════════════════════════════════════════════════════════════════

\subsection{Environment Setup}

The recommended workflow uses Conda to isolate the environment:

\begin{verbatim}
# 1. Clone repository
git clone https://github.com/YousefKad/Detecting-Breaks-Via-the-Fused-Lasso.git
cd Detecting-Breaks-Via-the-Fused-Lasso

# 2. Create and activate the environment
conda env create -f environment.yml
conda activate fused-lasso-breaks

# 3. Run the tests
pytest tests/ -v

# 4. Run the main Monte Carlo study
python -m simulations.N25_T5_m1
\end{verbatim}

\subsection{Key Dependencies}

\begin{center}
\begin{tabular}{lll}
  \toprule
  Package & Version & Purpose \\
  \midrule
  \texttt{numpy}      & $\ge 1.26$ & Array computation \\
  \texttt{cvxpy}      & $\ge 1.4$  & Convex optimisation \\
  \texttt{clarabel}   & $\ge 0.7$  & SOCP solver backend \\
  \texttt{matplotlib} & $\ge 3.8$  & Plotting \\
  \texttt{tabulate}   & $\ge 0.9$  & Table formatting \\
  \texttt{pytest}     & $\ge 8.0$  & Unit testing \\
  \bottomrule
\end{tabular}
\end{center}

\subsection{Reproducibility Note}

Monte Carlo results are stochastic by default (\texttt{SEED = None}).  To
reproduce a specific run exactly, set \texttt{SEED = <integer>} at the top of
the simulation file.  All numerical outputs are deterministic given a fixed
seed because CVXPY's CLARABEL solver is itself deterministic.

% ═════════════════════════════════════════════════════════════════════════════
\section{Connection to the Paper}
\label{sec:paper}
% ═════════════════════════════════════════════════════════════════════════════

\cref{tab:code_paper} maps the key objects in this codebase to their
counterparts in \citet{kaddoura2023estimation}.

\begin{table}[htb]
  \centering
  \caption{Correspondence between code and paper notation.}
  \label{tab:code_paper}
  \begin{tabular}{lll}
    \toprule
    Code symbol & Paper notation & Equation \\
    \midrule
    \texttt{y\_tilde} & $\tilde{y}_{it}$ & (2.3) \\
    \texttt{X\_tilde} & $\tilde{x}_{it}$ & (2.3) \\
    \texttt{b\_ols}   & $\hat{\bbeta}_t^{\text{OLS}}$ & (3.1) \\
    \texttt{lam}      & $\lambda$ & (3.2) \\
    \texttt{b\_hat}   & $\hat{\bB}_T(\lambda)$ & (3.2) \\
    \texttt{m\_hat}   & $\hat{m}(\lambda)$ & (3.3) \\
    \texttt{IC\_vec}  & $\text{IC}(\lambda)$ & (3.4) \\
    \texttt{lam\_star}& $\lambda^*$ & (3.5) \\
    \texttt{PHI}      & $\phi$   & Appendix A \\
    \texttt{PHI\_1}   & $\phi_1$ & Appendix A \\
    \texttt{PI}       & $\pi$    & Appendix A \\
    \bottomrule
  \end{tabular}
\end{table}

% ═════════════════════════════════════════════════════════════════════════════
\section*{References}
% ═════════════════════════════════════════════════════════════════════════════

\begin{thebibliography}{9}

\bibitem[Kaddoura and Westerlund(2023)]{kaddoura2023estimation}
Kaddoura, Y.\ and Westerlund, J.\ (2023).
\newblock Estimation of panel data models with random interactive effects and
  multiple structural breaks when $T$ is fixed.
\newblock \emph{Journal of Business \& Economic Statistics}, 41, 778--790.

\bibitem[Tibshirani et al.(2005)]{tibshirani2005}
Tibshirani, R., Saunders, M., Rosset, S., Zhu, J., and Knight, K.\ (2005).
\newblock Sparsity and smoothness via the fused lasso.
\newblock \emph{Journal of the Royal Statistical Society: Series B}, 67(1),
  91--108.

\bibitem[Zou(2006)]{zou2006}
Zou, H.\ (2006).
\newblock The adaptive lasso and its oracle properties.
\newblock \emph{Journal of the American Statistical Association}, 101(476),
  1418--1429.

\bibitem[Pesaran(2006)]{pesaran2006}
Pesaran, M.\ H.\ (2006).
\newblock Estimation and inference in large heterogeneous panels with a
  multifactor error structure.
\newblock \emph{Econometrica}, 74(4), 967--1012.

\bibitem[Bai(2010)]{bai2010}
Bai, J.\ (2010).
\newblock Common breaks in means and variances for panel data.
\newblock \emph{Journal of Econometrics}, 157(1), 78--92.

\end{thebibliography}

% ─────────────────────────────────────────────────────────────────────────────
\end{onehalfspacing}
\end{document}
